
\subsubsection{5.1 Main Result: Verbosity Coefficient Reversal (H-M2)}

The round-stratified logistic regression on HH-RLHF (53,600 rows per stratum, 2,000 bootstrap resamples) yields the following stylistic preference coefficients:

\textbf{Table 1: Stylistic preference coefficients across annotation strata (HH-RLHF)}

\begin{table}[htbp]
\centering
\begin{tabular}{lllllll}
\toprule
Feature & Early β & 95\% CI & Late β & 95\% CI & Δβ & Non-Overlapping CI \\
\midrule
β_L (verbosity) & −0.025 & [−0.043, −0.006] & +0.056 & [+0.043, +0.068] & **+0.080** & **Yes** \\
β_H (hedging) & −0.029 & [−0.048, −0.011] & −0.008 & [−0.024, +0.007] & +0.021 & No \\
β_S (structure) & −0.002 & [−0.021, +0.010] & +0.010 & [+0.004, +0.016] & +0.012 & No \\
β_Q (quality) & — & — & −0.017 & — & — & — \\
\bottomrule
\end{tabular}
\end{table}

Values are taken from \texttt{h-m2/experiment_results.json}. Model AUC: early stratum 0.495, late stratum 0.511.

Three observations:

\begin{enumerate}
\item \textbf{Verbosity preference reverses direction.} The 95\% bootstrap CIs for β_L are non-overlapping across strata. Early-stratum annotators are associated with a negative verbosity weight; late-stratum annotators are associated with a positive weight. The magnitude of the shift (Δβ_L = +0.080) is approximately 4× the early-stratum CI half-width.
\end{enumerate}

\begin{enumerate}
\item \textbf{Directional consistency across features, but only one reaches the non-overlap criterion.} All three stylistic Δβ values are positive (sign_consistent = true), but n_directional = 1 of 3, which is below the pre-registered threshold of 2. The H-M2 gate status is therefore PARTIAL, not PASS.
\end{enumerate}

\begin{enumerate}
\item \textbf{Quality covariate is stable.} The late-stratum β_Q = −0.017, well within the |β_Q| < 0.2 stability criterion, confirming that the verbosity shift does not reflect quality recalibration.
\end{enumerate}

\textbf{Note on model AUC.} AUC values near 0.5 reflect the high label noise inherent in large-scale pairwise preference annotation. The near-chance AUC does not invalidate coefficient-level inference: under n ≈ 53,600 per stratum, individual coefficient estimates and bootstrap confidence intervals remain statistically valid. The regression is designed to measure directionality of stylistic weighting, not to discriminate individual labels.

\textbf{Topic distribution imbalance.} A chi-square test for topic uniformity across strata yields p = 4.0×10⁻²⁷⁵, indicating extreme non-uniformity. This imbalance may confound the β_H and β_S estimates; response length (β_L) is considered less topic-sensitive and thus more robust to this confound.

\begin{figure}[htbp]
\centering
\includegraphics[width=\columnwidth]{figures/fig1_coefficient_comparison.png}
\caption{Coefficient comparison figure}
\label{fig:fig1_coefficient_comparison}
\end{figure}

\textit{Figure 1: Early-round and late-round logistic regression coefficients for verbosity (β_L), hedging (β_H), and structured reasoning (β_S) on HH-RLHF (53,600 rows per stratum, 2,000 bootstrap resamples). Error bars show 95\% CIs. β_L CIs are non-overlapping (Δβ_L = +0.080); β_H and β_S CIs overlap.}

\begin{figure}[htbp]
\centering
\includegraphics[width=\columnwidth]{figures/fig3_feature_stability_rounds.png}
\caption{Feature stability across rounds}
\label{fig:fig3_feature_stability_rounds}
\end{figure}

\textit{Figure 2: Trajectory of β_L, β_H, and β_S point estimates from stratum 1 to stratum 3. All three features trend positive, with β_L showing the largest and only non-overlapping CI shift.}

\subsubsection{5.2 AI-Typicality Geometric Projection (H-M1)}

The H-M1 experiment was designed to test within-annotator dose-response using PanelOLS fixed effects on WebGPT. This design could not be executed because worker identity fields are absent from the public WebGPT JSONL release. The analysis fell back to a between-group tercile design, as documented in \texttt{02c_experiment_brief.md} and confirmed in \texttt{h-m1/04_validation.md}.

Results from \texttt{h-m1/experiment_results.json}:

\begin{itemize}
\item Tercile F-statistic: \textbf{33.226} (p ≈ 8.3×10⁻⁹)
\item β_exposure (between-group regression): \textbf{4.1×10⁻⁵} (p ≈ 2.5×10⁻¹³)
\item Effect size criterion (≥ 0.1 SD per 1,000 tokens): \textbf{Not met} (\texttt{effect_size_ok = false})
\item Discriminant validity (placebo empirical p): \textbf{0.48} (confirming AI-typicality specificity)
\item HH-RLHF monotonicity: Not observed (\texttt{hh_monotonicity_ok = false})
\item n_webgpt_workers: 2 (tercile bin edges collapsed, limiting between-group power)
\end{itemize}

The significant tercile F-statistic indicates statistically significant between-group differences in AI-typicality projection scores. However, the pre-registered effect-size threshold was not met, and the planned within-annotator identification strategy was unavailable. The MUST_WORK gate for H-M1 is recorded as PASS on the basis that the code executed end-to-end, projection scores were computed, and the regression was estimable — not on the basis of scientific effect-size criteria.

The placebo discriminant validity result (empirical p = 0.48 for random unit vectors, vs. the observed significant tercile test) indicates that the between-group projection difference is specific to the AI-typicality direction rather than a general embedding artifact.

\begin{figure}[htbp]
\centering
\includegraphics[width=\columnwidth]{figures/dose_response.png}
\caption{Dose response figure}
\label{fig:dose_response}
\end{figure}

\textit{Figure 3: Between-group regression of AI-typicality projection score on exposure tercile (WebGPT, n=19,578). Tercile F=33.226, p≈8.3×10⁻⁹.}

\begin{figure}[htbp]
\centering
\includegraphics[width=\columnwidth]{figures/discriminant_validity.png}
\caption{Discriminant validity figure}
\label{fig:discriminant_validity}
\end{figure}

\textit{Figure 4: Discriminant validity — comparison of the observed between-group projection coefficient against the null distribution from 200 placebo permutations (random unit vectors). Empirical placebo p = 0.48, confirming signal specificity to the AI-typicality direction.}

\subsubsection{5.3 Existence Test and Null Results (H-E1)}

H-E1 tested whether directional stylistic coefficient drift exists in HH-RLHF after controlling for Q_early, and whether ambiguity modulates the drift. Results from \texttt{h-e1/experiment_results.json}:

\begin{itemize}
\item β_L Bonferroni-corrected p = 0.000 (nominally significant across rounds)
\item β_H Bonferroni p = 0.36; β_S Bonferroni p = 1.0
\item Round × ambiguity interaction p = 1.0; n_high_ambiguity = 0
\item VIF < 1.03 for all features (no multicollinearity)
\end{itemize}

The interaction test returned p = 1.0 because HH-RLHF lacks per-prompt disagreement labels; zero high-ambiguity samples were identified. This is a structural data limitation, not falsification of the interaction hypothesis. The H-E1 MUST_WORK gate is recorded as PASS.

\subsubsection{5.4 Cross-Experiment Summary}

\textbf{Table 2: Cross-experiment results}

\begin{table}[htbp]
\centering
\begin{tabular}{lllll}
\toprule
Sub-hypothesis & Gate Type & Gate Result & Primary Metric & Value \\
\midrule
H-E1 (existence) & MUST_WORK & PASS & β_L Bonferroni p & 0.000 \\
H-M1 (geometric) & MUST_WORK & PASS & Tercile F & 33.226, p≈8.3×10⁻⁹ \\
H-M2 (coefficient) & SHOULD_WORK & PARTIAL & n_directional & 1/3 \\
H-M3 (reward model) & SHOULD_WORK & NOT_STARTED & — & — \\
H-M4 (benchmark) & COULD_WORK & NOT_STARTED & — & — \\
\bottomrule
\end{tabular}
\end{table}

\textbf{Table 3: Informative null results and data requirements}

\begin{table}[htbp]
\centering
\begin{tabular}{lll}
\toprule
Null Result & Structural Reason & Data Required to Test \\
\midrule
Interaction p = 1.0 & No per-prompt disagreement labels in HH-RLHF & Multi-annotator Fleiss κ per prompt \\
Effect size below threshold & No within-annotator worker IDs in WebGPT public release & Genuine session tracking (≥3 sessions per worker) \\
β_H, β_S subthreshold & Index-based stratification + topic imbalance (p=4×10⁻²⁷⁵) & Topic-balanced strata, larger round depth \\
\bottomrule
\end{tabular}
\end{table}

\begin{figure}[htbp]
\centering
\includegraphics[width=\columnwidth]{figures/fig5_topic_balance.png}
\caption{Topic balance figure}
\label{fig:fig5_topic_balance}
\end{figure}

\textit{Figure 5: Chi-square test for topic uniformity across HH-RLHF annotation strata (p = 4×10⁻²⁷⁵). Extreme topic imbalance may confound β_H and β_S estimates but is less likely to affect β_L.}

